\chapter{Recognition and higher comparison}
\label{v4-ch:recognition-higher-comparison}

Two products can agree on cohomology while their primitives
produce different triple operations. The cubic example in the
prologue makes this distinction inside a strict differential
graded algebra: \(U^2=DH\), but
\(UH+HU=-V\) represents a nonzero cohomology class. Strict
associativity has encoded the higher operation; it has not
removed it. A comparison must transport the primitive as well as
the binary product.

Collision geometry imposes a further requirement. When the first
and third positions meet, the second input remains between them.
Moving an endpoint past that input uses an exchange map. Two
binary products can agree while these exchanges permit different
normal poles. We construct this example below, then give the
complexes and equations that decide whether a specified partial
comparison extends. The obstruction belongs to the operations
allowed by that comparison, with its boundary values fixed.

\section{Simultaneous operations and their recipients}
\label{v4-sec:recognition-operations}

Let \(k\) have characteristic zero and let \(X\) be the smooth
curve of Chapter~\ref{v4-ch:words-collisions-units}. For a finite
nonempty set \(I\), write
\(j_I:U_I\hookrightarrow X^I\) for the open immersion of the
locus where every pair of positions is distinct, and
\(\Delta_I:X\hookrightarrow X^I\) for the small diagonal.
For complexes of right differential modules \(M_i,N\) on \(X\),
define the complex of \emph{simultaneous chiral operations} by
\begin{equation}\label{v4-eq:recognition-chiral-operations}
 \mathcal Q_X(M_i;i\in I;N)
  =\operatorname{RHom}_{\mathcal D(X^I)}
       \left(j_{I,*}j_I^!\!\mathop{\boxtimes}_{i\in I}M_i,
                         \Delta_{I,*}N\right).
\end{equation}
The external tensors and geometric functors are derived. The
output retains the whole differential module on the diagonal,
including its normal derivatives. Degree-zero cycles represent
cochain operations in a chosen differential graded model; a
degree-minus-one cochain is a possible homotopy between such
operations. Passing to a scalar residue is a further map.
For ordinary modules, the degree-zero construction is the chiral
operation space of Beilinson and Drinfeld.\footnote{
A.~Beilinson and V.~Drinfeld, \emph{Chiral Algebras},
\S\S3.1.1--3.1.2, pp.~157--158. The composition below is the
derived form of their displayed composition of chiral operations.}

To compose, let \(\pi:I\twoheadrightarrow A\) have nonempty
fibres \(I_a\), so that \(I=\coprod_{a\in A}I_a\), and let
\(d_\pi:X^A\to X^I\) identify positions in each block.
Let \(j_\pi:U_\pi\hookrightarrow X^I\) be the open immersion
where positions in different blocks are distinct. Given operations
\[
 g_a\in\mathcal Q_X(M_i;i\in I_a;L_a),\qquad
 f\in\mathcal Q_X(L_a;a\in A;N),
\]
their composite is induced by
\begin{equation}\label{v4-eq:recognition-chiral-substitution}
 \begin{split}
 j_{I,*}j_I^!\!\mathop{\boxtimes}_I M_i
 &\simeq j_{\pi,*}j_\pi^!
       \mathop{\boxtimes}_A\left(j_{I_a,*}j_{I_a}^!
                                      \!\mathop{\boxtimes}_{I_a}M_i\right)\\
 &\xrightarrow{\,\boxtimes g_a\,}
       j_{\pi,*}j_\pi^!\!\mathop{\boxtimes}_A\Delta_{I_a,*}L_a\\
 &\simeq d_{\pi,*}j_{A,*}j_A^!\!\mathop{\boxtimes}_A L_a
       \xrightarrow{\,d_{\pi,*}f\,}\Delta_{I,*}N.
 \end{split}
\end{equation}
The middle equivalence is open--closed base change. Every
parenthesization of a nested partition uses the same disjoint
open set and the same composite diagonal map. Coherent base
change therefore makes these compositions associative, with all
their homotopies. The unary identity is the identity of a
differential module. Permuting labels transports the same
diagram. Thus the operation complexes form a coherent coloured operad:
the colours are the chosen differential modules, and an
operation has its indicated list of input colours and output
colour. Forgetting permutations leaves its nonsymmetric operad
on ordered inputs. For equations between cochains, choose a
differential graded operad model compatible with these geometric
functors and their base-change maps.

The word carrier has a different endomorphism operad. For word
systems \(F_\lambda\in\mathcal W_X\), its complexes are
\begin{equation}\label{v4-eq:recognition-word-operations}
 \mathcal P_F(\lambda_1,\ldots,\lambda_r;\nu)
   =\operatorname{RHom}_{\mathcal W_X}
       (F_{\lambda_1}\star\cdots\star F_{\lambda_r},F_\nu),
 \qquad r\geq0,
\end{equation}
where \(\star=\star_{\mathrm{cut}}\) and the empty tensor is
\(K\). Substitution is
\(f\circ(g_1\star\cdots\star g_r)\), with the coherent cut
associator. On homogeneous blocks \(x_i\), the tensor of maps
has sign \((-1)^{\sum_{i<j}|g_j||x_i|}\). The cut construction
of Chapter~\ref{v4-ch:words-collisions-units} proves its operad
coherences: both composites retain the same final consecutive
decomposition and localized coefficient. Its arity-zero maps
have source \(K\), not the local vacuum \(J\).

A field comparison has to relate independently specified
operations to \eqref{v4-eq:recognition-word-operations}. For a
chosen family whose operations are represented by a differential
graded operad \(\mathcal Q^{\mathrm{adm}}\), it includes maps
\begin{equation}\label{v4-eq:recognition-comparison}
 \rho_r:\mathcal Q^{\mathrm{adm}}
                 (\lambda_1,\ldots,\lambda_r;\nu)
       \longrightarrow
                 \mathcal P_F(\lambda_1,\ldots,\lambda_r;\nu)
\end{equation}
compatible with differentials and the required substitutions.
Here \emph{admissible} means a specified complex of operations, with
its coefficient domains, pole conditions and topology. In a
meromorphic differential-module family it can be a subcomplex
of \eqref{v4-eq:recognition-chiral-operations}; closure under the
required substitutions must be proved. A fixed pole bound need
not have this property, since compositions can increase pole
order. Monodromic or analytic fields require their own coefficient
recipient before such a map is defined.

Collision restrictions, exchanges, translation and local-vacuum
insertion give additional diagrams for \(\rho\). When their
commutativity is only up to homotopy, the homotopies and their
higher compatibilities are part of the comparison. The
construction of either operad above does not produce
\eqref{v4-eq:recognition-comparison}. Nor does the singleton
\(F_1\) determine its target, which uses the whole word system.

\begin{example}[An undetected positive component]
\label{v4-ex:recognition-singleton}
For a differential module \(M\), let
\(I(M)_n=\Delta_{n,*}M\) for \(n>0\), with zero empty
component. Closed base change gives its monotone collision
equivalences. Put \(I=I(J_X)\), where \(J_X=p_X^!k\), and
\(H=I\star I\). Then \(H_1=0\), while \(H_2\) is nonzero
on the separated two-point locus. The free augmented cut
algebras
\[
 \mathbb T(I)\longrightarrow\mathbb T(I\oplus H),
 \qquad \mathbb T(E)=K\oplus\bigoplus_{r\geq1}E^{\star r},
\]
agree at the singleton and differ at two positions by the extra
summand \(H_2\). Each sum is finite in every geometric arity.
Consequently singleton evaluation is not conservative even on
connected augmented cut algebras.
\end{example}

\section{The endpoint exchange at a ternary collision}
\label{v4-sec:recognition-spectator}

Let \(B=k[\epsilon]/(\epsilon^2)\), concentrated in
cohomological degree zero. Give it an auxiliary parity by making
\(1\) even and \(\epsilon\) odd. For \(\sigma=\pm1\), define
\[
 R_\sigma(a\otimes b)=\sigma^{p(a)p(b)}b\otimes a
\]
on the two homogeneous basis vectors, and extend linearly.
The parity here specifies the exchange; it is not a
cohomological Koszul sign. Both exchanges have the same unit
crossings, are involutive, and satisfy the braid relation.
For three basis vectors the two braid words have the same sign
\(\sigma^{p(a)p(b)+p(a)p(c)+p(b)p(c)}\).
They also satisfy the two multiplication compatibility identities
\[
 \begin{split}
 R_\sigma(1\otimes\mu)
   &=(\mu\otimes1)(1\otimes R_\sigma)(R_\sigma\otimes1),\\
 R_\sigma(\mu\otimes1)
   &=(1\otimes\mu)(R_\sigma\otimes1)(1\otimes R_\sigma).
 \end{split}
\]
When the multiplied entries are not both \(\epsilon\), their
parities add; when both are \(\epsilon\), both sides vanish.
This proves the identities on a basis. The binary multiplication
\(\mu\) is the same for both signs, and \(\mu R_\sigma=\mu\).

To collide the first and third input while retaining the middle
input, define
\begin{equation}\label{v4-eq:recognition-outer-quotient}
 q_\sigma=(\mu\otimes1)(1\otimes R_\sigma):
              B^{\otimes3}\longrightarrow B\otimes B.
\end{equation}
The first output is the product of the endpoints; the second is
the unchanged middle state after exchange. For
\(\xi=\epsilon\otimes\epsilon\otimes1
          -1\otimes\epsilon\otimes\epsilon\), one obtains
\begin{equation}\label{v4-eq:recognition-spectator-witness}
 q_+(\xi)=0,\qquad q_-(\xi)=2\epsilon\otimes\epsilon.
\end{equation}
Thus an operation that vanishes for one exchange need not vanish
for the other, although the binary products agree.

This discrepancy gives actual differential modules with different
allowed poles. On the affine-line three-point chart put
\[
 Y=\operatorname{Spec}
 k[z_1,z_2,z_3,(z_1-z_2)^{-1},(z_2-z_3)^{-1}],
 \qquad u=z_1-z_3,\qquad T=B^{\otimes3}.
\]
Inside \(T\otimes_k\omega_Y[u^{-1}]\), define right
differential modules
\begin{equation}\label{v4-eq:recognition-polar-extensions}
 \mathcal M_\sigma
   =T\otimes_k\omega_Y
        +(\ker q_\sigma)\otimes_k\omega_Y[u^{-1}].
\end{equation}
The kernels are constant coefficient spaces, so multiplication
by regular functions and differentiation preserve these modules.
In a coordinate density, the right derivative differentiates
the coefficient with a minus sign. Every finite normal
derivative of an allowed pole remains allowed.

\begin{proposition}[Binary agreement and a ternary failure]
\label{v4-prop:recognition-ternary-poles}
Both modules in \eqref{v4-eq:recognition-polar-extensions} restrict
to the same meromorphic module off \(u=0\). The identity of that
module extends to no map \(\mathcal M_+\to\mathcal M_-\), and
to no map in the opposite direction. Their common binary
coefficient extension, obtained from \(\ker\mu\), does not
decide this ternary extension problem.
\end{proposition}

\begin{proof}
Localizing at \(u\) makes the first summand in
\eqref{v4-eq:recognition-polar-extensions} the whole meromorphic
module. The target is a submodule of that localization and has
no \(u\)-torsion. A map restricting to the identity must therefore
be the identity on every section: the difference vanishes after
localization and hence vanishes already.

But \(u^{-1}\xi\,dz_1dz_2dz_3\) belongs to \(\mathcal M_+\)
and not to \(\mathcal M_-\), by
\eqref{v4-eq:recognition-spectator-witness}. A nonzero constant
coefficient of \(u^{-1}\) cannot become regular at \(u=0\).
For the reverse direction, use
\(\eta=\epsilon\otimes\epsilon\otimes1
          +1\otimes\epsilon\otimes\epsilon\):
\(q_-(\eta)=0\) and \(q_+(\eta)=2\epsilon\otimes\epsilon\).
At two inputs the construction uses the common map
\(\mu:B^{\otimes2}\to B\), so its kernel extension is identical.
\end{proof}

The proposition fixes the identity on the marked meromorphic
coefficient space. It makes no assertion about isomorphisms that
change those markings. It also constructs binary and ternary
coefficient extensions, not a complete field system in every
arity. An all-arity realization must additionally supply the
collision maps and verify every substitution among them.

\section{The complex controlling a higher comparison}
\label{v4-sec:recognition-comparison-tower}

Let \(R\) be a commutative ring. We use the suspended convention
of Chapter~\ref{v4-ch:bar-cobar-topology}. An augmented
\(A_\infty\)-algebra is specified here by a complex
\(V=s\bar A\) and a degree-one coderivation
\[
 b_A:T_R^c(V)\longrightarrow T_R^c(V),\qquad b_A^2=0,
 \qquad b_A(1)=0.
\]
Its Taylor components are maps \(b_{A,r}:V^{\otimes r}\to V\)
of degree one, with \(b_{A,1}\) the differential on \(V\).
The augmented state algebra is the unitization
\(R\oplus V[-1]\): its new unit acts by the identity under
binary multiplication, and every higher operation vanishes when
one of its inputs is that unit. All displayed inputs belong to
the augmentation ideal. For a differential graded
algebra, only \(b_1,b_2\) occur, with
\(b_1(sa)=-s(da)\) and \(b_2(sa,sb)=(-1)^{|a|}s(ab)\).

For a second such algebra put \(W=s\bar B\). A normalized
\(A_\infty\)-map is a coaugmentation-preserving degree-zero
coalgebra map
\(F:T_R^c(V)\to T_R^c(W)\) satisfying
\(b_BF=Fb_A\). Its components
\(f_r:V^{\otimes r}\to W\) determine it by consecutive
subdivisions. Define the actual cochain complex
\begin{equation}\label{v4-eq:recognition-hom-complex}
 E_r=\operatorname{Hom}_R(V^{\otimes r},W),\qquad
 \partial f=b_{B,1}f-(-1)^{|f|}f\,b_{A,1}^{(r)},
\end{equation}
where \(b_{A,1}^{(r)}\) is the tensor differential. For derived
comparisons use specified cofibrant coefficient models, or
replace this Hom by its derived model. The formulas below concern
the chosen chain models, with no unmentioned resolution change.

Suppose \(f_1,\ldots,f_{n-1}\) already satisfy the comparison
equations below length \(n\). Define
\begin{equation}\label{v4-eq:recognition-obstruction}
 \begin{split}
 O_n={}&
 \sum_{\substack{q\geq2,\ i_1+\cdots+i_q=n\\i_j\geq1}}
       b_{B,q}(f_{i_1}\otimes\cdots\otimes f_{i_q})\\
 &-\sum_{\substack{u+v+w=n\\v\geq2}}
       f_{u+1+w}(1^{\otimes u}\otimes b_{A,v}\otimes1^{\otimes w})
       \in E_n^1.
 \end{split}
\end{equation}
All components on the right have index less than \(n\), except
the already specified algebra operations. The signs are tensor
composition signs: insertion of \(b_{A,v}\) after a homogeneous
prefix has the sign given by its total degree. No division by
two occurs.

\begin{theorem}[One higher component]
\label{v4-thm:recognition-one-step}
The cochain \(O_n\) is closed. A component \(f_n\) extending
the comparison through length \(n\) exists precisely when
\([O_n]=0\) in \(H^1(E_n)\). It must satisfy
\[
 \partial f_n=-O_n.
\]
When a solution exists, all solutions form a torsor for
\(Z^0(E_n)\). Modulo changes \(f_n\mapsto f_n+\partial h\),
with \(h\in E_n^{-1}\) and all lower components fixed, they
form a torsor for \(H^0(E_n)\).
\end{theorem}

\begin{proof}
Extend the chosen components to a coalgebra map by setting the
remaining Taylor components to zero. Put \(D=b_BF-Fb_A\).
The square-zero equations imply
\[
 b_BD+Db_A=0.
\]
The map \(D\) is a coderivation over \(F\). Since its
length-one output vanishes on inputs shorter than \(n\), the
only part of this identity with \(n\) inputs and one output is
\(b_{B,1}O_n+O_n b_{A,1}^{(n)}=0\). This is
\(\partial O_n=0\). In the equation with one output and
\(n\) inputs, the new \(f_n\) occurs only in the unary terms;
all other terms are \eqref{v4-eq:recognition-obstruction}.
This gives the displayed equation. Subtracting two solutions
gives a degree-zero cocycle, and quotienting by degree-zero
boundaries gives the final assertion.
\end{proof}

For strict differential graded inputs, the first equation is
\[
 \partial f_2
   =f_1b_{A,2}-b_{B,2}(f_1\otimes f_1).
\]
Thus \(f_2\) is the multiplicativity homotopy in suspended
coordinates. At length three, \(f_2\) enters on both sides of
the next equation. Choosing its cohomology class independently
of that equation can obstruct the extension.

\begin{example}[A nonzero first ternary obstruction]
\label{v4-ex:recognition-minimal-ternary}
Let \(R=k\), and take \(V=kx\oplus ky\) with
\(|x|=-1\), \(|y|=-2\). On \(T^c(V)\), let the first
coderivation be zero and let the second have the sole component
\[
 b_{B,3}(x,x,x)=y,
\]
vanishing on inputs containing \(y\). Its square is zero:
nested insertions vanish, while disjoint degree-one insertions
cancel. With \(f_1=1_V\), both unary and binary comparison
equations hold, for any \(f_2\). But \(O_3(x,x,x)=y\),
and \(\partial=0\) on \(E_3\). Hence no higher comparison
with this linear term exists. The corresponding unsuspended
states have degrees zero and minus one.
\end{example}

There is a version on the actual word carrier. Take the positive
system \(I=I(J_X)\) of
Example~\ref{v4-ex:recognition-singleton}, put \(H=I^{\star3}\),
and use reduced states \(G=I\oplus H[1]\). On the suspended
states \(G[1]=I[1]\oplus H[2]\), define \(b_3\) to be the
degree-one shifted identity
\[
 (I[1])^{\star3}\simeq H[3]\longrightarrow H[2],
\]
and zero if an input belongs to \(H[2]\). Together with the
internal differential it squares to zero by the same cancellation.
Compare it to the structure with \(b_3=0\) and the same unary
differential and zero binary product. The first ternary defect
projects to the identity of \(H\). It is not a boundary, because
\(H_3\) is nonzero on the separated three-point locus. Thus
the identity of the underlying word system and the agreement of
binary operations do not extend to an \(A_\infty\)-map. This
example uses positive systems and lies within the geometric
finiteness regime; coefficient completion is not its obstruction.

\section{Fixed boundary values and allowed homotopies}
\label{v4-sec:recognition-relative-complex}

Geometric conditions restrict the possible \(f_n\). A
restriction to an open set, a prescribed normal coefficient or
a vacuum value can be represented, on fixed chain models, by a
cochain map
\(r_n:E_n\to T_n\). Let \(t_n\in T_n^0\) be the
prescribed value of the new component. Compatibility with its
chain equation requires
\[
 \partial t_n=-r_nO_n.
\]
Several restrictions can be placed in the same target diagram;
its maps and compatibility homotopies must be included when
forming the derived restriction map.

First suppose \(r_n\) is degreewise surjective and require the
strict boundary value \(r_nf_n=t_n\). Put \(K_n=\ker r_n\),
and choose any degree-zero lift \(\widetilde f_n\) of \(t_n\).

\begin{theorem}[The relative obstruction]
\label{v4-thm:recognition-relative}
The cochain
\[
 o_n=O_n+\partial\widetilde f_n
\]
is a degree-one cocycle in \(K_n\). Its class in \(H^1(K_n)\)
is independent of the lift. The comparison extends with the
prescribed boundary value if and only if this class vanishes.
When it does, solutions modulo homotopies in \(K_n^{-1}\),
with all lower data fixed, form a torsor for \(H^0(K_n)\).
\end{theorem}

\begin{proof}
The boundary equation gives \(r_no_n=0\), and
\(\partial O_n=0\) gives \(\partial o_n=0\). Changing the
lift by \(a\in K_n^0\) changes \(o_n\) by \(\partial a\).
A correction \(a\in K_n^0\) makes
\(f_n=\widetilde f_n+a\) a solution precisely when
\(\partial a=-o_n\). Subtraction of two solutions and then
quotienting by relative boundaries gives the asserted torsor.
\end{proof}

Without surjectivity, a strict lift of \(t_n\) is a separate
first requirement. If the boundary is instead allowed to agree
up to a specified homotopy, the correct complex is the
\emph{homotopy fibre}
\begin{equation}\label{v4-eq:recognition-fibre}
 \operatorname{Fib}(r_n)^q=E_n^q\oplus T_n^{q-1},
 \qquad d(e,h)=(\partial e,r_ne-\partial h).
\end{equation}
The pair \((O_n,-t_n)\) is closed. Its class vanishes precisely
when there are \(f_n\in E_n^0\), \(h_n\in T_n^{-1}\) with
\[
 \partial f_n=-O_n,\qquad
 r_nf_n=t_n-\partial h_n.
\]
Indeed, take the negative of the \(E_n\)-component of a
primitive of \((O_n,-t_n)\), and retain its \(T_n[-1]\)
component. This proves the equivalence directly, including the
sign of the boundary homotopy. Replacing a strict boundary
condition by this homotopy condition changes the comparison
problem unless that replacement was allowed.

\begin{example}[An inadmissible primitive]
\label{v4-ex:recognition-forbidden-primitive}
Let \(E\) have basis \(h\) in degree zero and \(c\) in
degree one, with \(\partial h=c\). Let \(T=kh\) in degree
zero and let \(r:E\to T\) be the quotient by \(kc\).
The obstruction \(O=c\) is exact in \(E\). With prescribed
boundary \(t=0\), however, \(K=kc\) has
\(H^1(K)=kc\), so no allowed primitive exists. The unique
ambient solution \(-h\) has the wrong boundary value. This is
the distinction between finding a primitive and finding one in
the permitted support or vacuum sector.
\end{example}

\begin{corollary}[Extension through all lengths]
\label{v4-cor:recognition-all-lengths}
Suppose the restriction complexes are specified at every length,
each \(r_n\) is degreewise surjective, and \(H^1(K_n)=0\)
for \(n\geq2\). An admissible unary comparison extends to a
normalized \(A_\infty\)-map if compatible boundary values can
be supplied successively for its partial lifts. If contracting homotopies
\(h_n:K_n\to K_n[-1]\), satisfying
\(\partial h_n+h_n\partial=1\), are supplied, one can take
\[
 f_n=\widetilde f_n-h_no_n.
\]
Over a field, a quasi-isomorphism \(f_1\) gives a
quasi-isomorphism of the resulting direct-sum bar complexes.
\end{corollary}

\begin{proof}
Apply Theorem~\ref{v4-thm:recognition-relative} recursively. For
the displayed choice, closedness of \(o_n\) gives
\(\partial h_no_n=o_n\). At each finite word length the
coalgebra formula uses only finitely many components, so the
resulting map is defined on the direct-sum bars and satisfies
the chain identity there.

Filter these bars by word length. On the associated graded, the
comparison is \(f_1^{\otimes r}\) in length \(r\). Over a
field tensor powers preserve quasi-isomorphisms. Finite-stage
induction and exact filtered unions prove the final assertion.
\end{proof}

The hypotheses include compatibility of the prescribed boundary
values with the partial solution chosen at the previous stage.
Vanishing cohomology does not supply incompatible boundary data.
Nor does it by itself imply uniqueness in a full mapping space:
\(H^0(K_n)\) describes this stage's choices, and negative-degree
cohomology can describe homotopies between those choices. For
completed coefficient bars, the same construction must also
satisfy the continuity conditions of
Chapter~\ref{v4-ch:bar-cobar-topology}.

\section{Operations, deformations, and a recognition problem}
\label{v4-sec:recognition-deformations}

There is a parallel obstruction for constructing the operations
themselves. Fix \(b_1\) on a suspended coefficient complex \(V\)
and choose \(b_2,\ldots,b_{n-1}\) satisfying \(b^2=0\) below
length \(n\). Then
\[
 C_n=\sum_{\substack{i+j=n+1\\i,j\geq2}}
       \sum_{u=0}^{i-1}
       b_i(1^{\otimes u}\otimes b_j\otimes1^{\otimes(i-1-u)})
       \in\operatorname{Hom}_R(V^{\otimes n},V)^2.
\]
Its differential is zero. To see this without dividing by any
integer, use the identity \(b\,b^2=b^2b\) for the partial
coderivation. Its first nonzero component is
\(b_1C_n-C_nb_1^{(n)}=0\). A new operation \(b_n\) must
satisfy \(\partial b_n=-C_n\), so the obstruction is in
degree two and the choices, modulo boundaries, are in degree
one. This differs by one degree from the comparison problem,
where \(f_n\) has degree zero.

In characteristic zero these equations also have a Lie-algebra
form. Let
\[
 \mathfrak g=\prod_{r\geq2}
                   \operatorname{Hom}_k(V^{\otimes r},V),
 \qquad d_{\mathfrak g}=[b_1,-].
\]
Identify a collection with its coderivation of \(T^c(V)\).
The bracket is the graded commutator of coderivations, projected
back to one output. Give arity \(r\) weight \(r-1\), and
let \(F^p\mathfrak g\) consist of collections supported in
weights at least \(p\). Then
\([F^p,F^q]\subset F^{p+q}\), and the product is complete
and separated for this filtration. Each fixed-arity bracket
coefficient is a finite sum.

A degree-one \(\alpha\in\mathfrak g\) gives square-zero
\(b_1+\alpha\) precisely when
\[
 d_{\mathfrak g}\alpha+\tfrac12[\alpha,\alpha]=0.
\]
If \(\alpha\) solves this equation modulo \(F^p\), its
curvature in \(F^p/F^{p+1}\) is a degree-two cocycle.
The identity
\(d_{\mathfrak g}C+[\alpha,C]=0\) follows by expanding the
curvature, and its leading term proves closedness. A correction
\(\beta\in F^p\mathfrak g^1\) changes that class by
\(d_{\mathfrak g}\beta\). This gives the same extension
criterion as the integral coderivation equation. For
deformations around a complete nonzero solution, the controlling
differential is instead
\(d_\alpha=d_{\mathfrak g}+[\alpha,-]\); retaining that twist
is essential. These formulas apply to an admissible operation
subcomplex only if it is closed under the differential and the
required insertions.

To compare an entire coloured operation system, one can specify
a differential graded operad resolution \(\mathcal R\) with
successive free generators \(e\), each having differential in
the previously constructed part. If a partial map
\(\rho:\mathcal R\to\mathcal P_F\) is fixed, extending it
over a generator of degree \(d\) asks for
\[
 x\in\mathcal P_F^d,\qquad dx=\rho(de).
\]
The right side is closed because \(d^2e=0\). Its class in
\(H^{d+1}(\mathcal P_F)\), or in the relative complex with
the prescribed boundary values, is exactly the obstruction.
Freeness then extends the map to composites; the chain identity
on generators implies it on every composite by the Leibniz
rule. This is an explicit obstruction calculation once the
generators, their differentials and their operation recipients
are supplied. A name for an operad or a Hochschild group supplies
none of those data.

\begin{question}[Recognition of an arithmetic family]
\label{v4-qn:recognition-arithmetic}
For an independently specified arithmetic coefficient family,
construct its simultaneous operation complexes and their
geometric recipients, then the maps
\eqref{v4-eq:recognition-comparison} on their full coefficient
domains. Compare every differential, supported restriction,
normal derivative, endpoint exchange and local-vacuum insertion
used by the family, with compatible nested substitutions.
Retain arithmetic operators and a state or trace when these
marks have been independently defined. Determine the image of
the resulting comparison, or prove an equivalence on a stated
class of objects and full mapping spaces.
\end{question}

The preceding constructions locate this question's successive
equations and give an extension theorem when their relative
cohomology vanishes. The ternary pole example shows why the
binary field product cannot replace the simultaneous operation
system. The cubic example shows why a strict presentation cannot
replace the comparison of primitives. Both distinctions persist
when a scalar character or trace already agrees.
